\section{The RS04 codec}
\label{rs04}
\sloppy

This section describes the data format of the dvdisaster RS04 codec.
{\bf Status: draft.} dvdisaster implements RS04 as an experimental codec
(command line only); the format may still change until this section is
declared final (see the changelog).

\medskip

RS04 keeps the basic idea of RS03 -- the medium is divided into layers and
error correction block (``codeword'') $i$ consists of sector $i$ of every
layer -- but changes almost everything else:

\begin{itemize}
\item The Reed-Solomon code works over GF($2^{16}$) instead of GF($2^8$).
  A codeword has up to 65535 symbols instead of 255, so the number of layers
  $n$ can be much larger and each codeword is spread over many more places on
  the medium. Against scattered damage this brings the tolerated loss close
  to the theoretical limit, while damage at the outer edge is handled as well
  as by RS03 (section \ref{rs04-rationale}).
\item The layout is {\em self-describing}. Small {\em tag sectors} are
  replicated throughout the error correction data. Each of them describes the
  complete layout, so a single readable tag is enough. No table of medium
  sizes and no brute-force search are needed, and the augmented image may
  have any size.
\item Several RS04 sets may exist on one medium (one per session of a
  multi-session disc). All positions are relative to the start of the set.
\item Besides 2048-byte sectors, Mode 2 Form 2 sectors (2324 bytes) and
  CD-DA audio sectors (2352 bytes) can be protected. Audio extraction with
  a different drive read offset is handled by an offset search.
\item The error correction file (``ecc file'') is byte-for-byte the same
  data that is appended to the image in the augmented case.
\end{itemize}

Everything in this section is normative unless it is marked as
{\em recommended} or {\em informative}. Recommended procedures describe
what dvdisaster does; other encoders may choose differently, and decoders
must accept every set which satisfies the validity rules.

\subsection{Conventions}
\label{rs04-conventions}

\begin{itemize}
\item A {\em sector} is a block of $B$ bytes, where $B$ is the sector size of
  the set (section \ref{rs04-formats}). Positions and counts are given in
  sectors and start at 0.
\item All integers on disc are stored {\em little-endian}, independent of the
  host. Big-endian hosts must convert them after reading and before writing
  (as dvdisaster does for RS03 with {\tt SwapEccHeaderBytes()}).
\item $\lceil x \rceil$ and $\lfloor x \rfloor$ round up and down.
  ${\rm NextPow2}(x)$ is the smallest power of two $\ge x$
  (${\rm NextPow2}(1)=1$), ${\rm Pow2Floor}(x)$ the largest power of two
  $\le x$ for $x \ge 1$, and ${\rm Pow2Floor}(0)=1$.
\item {\em Crc32} is the function {\tt Crc32()} of dvdisaster which is also
  used by RS01--RS03: reflected polynomial {\tt 0xEDB88320}, initial value
  {\tt 0xFFFFFFFF}, {\em no} final inversion. It equals the zlib
  {\tt crc32()} value XOR {\tt 0xFFFFFFFF}. Check values:
  Crc32(``123456789'') = {\tt 0x340BC6D9}; 2048 zero bytes give
  {\tt 0x0E174561}, 2324 zero bytes {\tt 0x2F6ED567} and 2352 zero bytes
  {\tt 0x416831C0}. Crc32 values are stored as little-endian 32-bit integers.
\item {\em MD5} is the message digest of RFC 1321.
\end{itemize}

\subsection{Terminology}
\label{rs04-terms}

An {\em RS04 set} consists of the protected {\em data sectors} and the
{\em added zone} which RS04 creates. The added zone holds tag sectors, CRC
sectors and parity sectors. The following symbols are used throughout this
section:

\medskip
\begin{tabular}{|l|p{12.5cm}|}
\hline
$B$ & sector size: 2048, 2324 or 2352 bytes \\
$B_s$ & shard size $64 \lceil B/64 \rceil$: 2048 or 2368 bytes \\
$S$ & number of data sectors, $S = \lceil dataBytes / B \rceil \ge 1$ \\
$c$ & number of CRC sectors, $c = \lceil S / 512 \rceil$ \\
$D$ & number of protected sectors, $D = S + c$ \\
$L$ & layer size (sectors per layer) $\ge 1$ \\
$k$ & number of data layers = original symbols per codeword; $kL \ge D$ \\
$m$ & number of ecc layers = parity symbols per codeword; $1 \le m \le k$ \\
$n$ & codeword length $n = k + m$ \\
$M$ & ${\rm NextPow2}(m)$; the code requires $M + k \le 65536$ \\
$T$ & tag interval: one tag sector every $T$ sectors of the added zone \\
$Q$ & number of CRC and parity sectors, $Q = c + mL$ \\
$A$ & size of the added zone, $A = Q + \lceil Q / (T-1) \rceil$ \\
\hline
\end{tabular}
\medskip

The stored parameters are $B$, $dataBytes$, $L$, $k$, $m$ and $T$; all other
values are derived from them.

{\footnotesize
\begin{verbatim}
relative position 0                  S                           S+A
|<--------- data sectors -------->|<------- added zone ------->|
| D_0 D_1 ...               D_S-1 | t C..C P..P t P..P t P..P  |
   t = tag (every T sectors), C = CRC sector, P = parity sector

logical view       column 0    column 1   ...  column L-1
  data layer 0     X_0         X_1             X_L-1
  data layer 1     X_L         X_L+1           X_2L-1
  ...
  data layer k-1   X_(k-1)L    ...             X_kL-1
  ecc layer 0      P_0         P_1             P_L-1
  ...
  ecc layer m-1    P_(m-1)L    ...             P_mL-1

X_t = D_t (t < S), C_(t-S) (S <= t < D), zero sector (D <= t < kL)
codeword i = column i
\end{verbatim}
}

\subsection{Sector formats}
\label{rs04-formats}

\begin{tabular}{|r|p{7.6cm}|r|r|}
\hline
$B$ & Contents of a data sector & $B_s$ & tail $r$\\
\hline
2048 & user data of Mode 1 and Mode 2 Form 1 CD sectors, DVD and BD
       sectors (ISO and UDF images) & 2048 & 0 \\
2324 & user data of Mode 2 Form 2 CD sectors (e.g. VCD/SVCD video) & 2368 & 20\\
2352 & CD-DA audio: 588 stereo samples of 4 bytes each (16-bit signed
       little-endian left channel, then right channel), as returned by
       READ CD and stored in {\tt .bin} rips & 2368 & 48 \\
\hline
\end{tabular}
\medskip

All sectors of a set -- data, tag, CRC and parity sectors -- have exactly
$B$ bytes. Augmented images and ecc files store them back-to-back without
any per-sector framing (no sync pattern, header, subheader or EDC/ECC of
the physical CD sector). The sector size is stored in every tag.

\paragraph{Last data sector.} The data is $dataBytes$ bytes long. If
$dataBytes$ is not a multiple of $B$, the last data sector is padded with
zero bytes to $B$ bytes for encoding and for its Crc32. The padding is not
part of the data; it is not covered by the MD5 sum. Augmented images
require $dataBytes = S \cdot B$, otherwise the added zone would not start on
a sector boundary. Ecc files may protect data of any length, e.g. arbitrary
files.

\paragraph{Missing sectors in image files.}
\label{rs04-missing}
dvdisaster stores a sector which could not be read as its {\em dead sector
marker}: 2048 bytes starting with the text {\tt dvdisaster dead sector
marker} (the marker used for RS01--RS03 images), or 2048 copies of a fill
byte if the user requested one (e.g.\ zero for images made by other
tools). For $B > 2048$ a missing sector is this 2048-byte marker followed by
$B - 2048$ zero bytes (fill bytes if a fill byte is used). Readers recognise
a missing sector by its first 2048 bytes only. A missing sector is an
erasure; its Crc32 is meaningless.

\paragraph{Whole-disc images (informative).} dvdisaster reads a
multi-session medium as one image in which every sector is at its absolute
position. The sectors between the sessions (lead-out, lead-in and pregap) can
not be read; they are stored as zero-filled $B$-byte sectors. They belong to
no set (section \ref{rs04-sessions}) and are not erasures.

\subsubsection{Sector to shard mapping}
\label{rs04-shard}

The Reed-Solomon code (section \ref{rs04-code}) works on {\em shards} of
$B_s$ bytes in the buffer format of Leopard-RS: the shard is a sequence of
64-byte blocks, and byte $u$ ($0 \le u < 32$) of block $b$ is the low byte,
byte $u+32$ the high byte of the 16-bit symbol $32b+u$. A sector is turned
into a shard as follows, with $F = 64 \lfloor B/64 \rfloor$ and $r = B - F$:

\begin{itemize}
\item shard bytes $0 \ldots F-1$ = sector bytes $0 \ldots F-1$;
\item shard bytes $F \ldots F + r/2 - 1$ = sector bytes $F \ldots F + r/2 - 1$
  (low bytes of the first $r/2$ symbols of the last block);
\item shard bytes $F+32 \ldots F + 32 + r/2 - 1$ = sector bytes
  $F + r/2 \ldots B-1$ (high bytes of the same symbols);
\item all other shard bytes are zero.
\end{itemize}

For $B = 2048$ the shard is the sector itself. For a parity shard the bytes
which are zero in this mapping are always zero: the symbols they belong to
are zero in every original shard and the code is linear. Therefore a parity
shard is turned back into a $B$-byte parity sector by the inverse mapping,
and no padding is ever stored. Decoders may check that the unused shard
bytes of decoded shards are zero.

{\em Note:} simply appending zero bytes to the sector would not work for
$B \ne 2048$, since the parity would then have non-zero high bytes for the
partially used symbols of the last block (20 resp.\ 24 extra bytes) which do
not fit into a $B$-byte parity sector.

\subsection{The Reed-Solomon code}
\label{rs04-code}

\subsubsection{Field and symbols}

RS04 computes in GF($2^{16}$) = GF(2)$[z] / (p(z))$ with
$p(z) = z^{16} + z^5 + z^3 + z^2 + 1$ ({\tt 0x1002D}). Writing a field
element as a 16-bit integer whose bit $b$ is the coefficient of $z^b$ gives
its {\em polynomial representation}. The symbols stored on disc use a
different representation, based on the following {\em Cantor basis}
$\beta_0, \ldots, \beta_{15}$ (given in polynomial representation):

\medskip
\begin{center}
\begin{tabular}{|c|c|c|c|c|c|c|c|}
\hline
$\beta_0$ & $\beta_1$ & $\beta_2$ & $\beta_3$ & $\beta_4$ & $\beta_5$ & $\beta_6$ & $\beta_7$ \\
\tt 0001 & \tt ACCA & \tt 3C0E & \tt 163E & \tt C582 & \tt ED2E & \tt 914C & \tt 4012 \\
\hline
$\beta_8$ & $\beta_9$ & $\beta_{10}$ & $\beta_{11}$ & $\beta_{12}$ & $\beta_{13}$ & $\beta_{14}$ & $\beta_{15}$ \\
\tt 6C98 & \tt 10D8 & \tt 6A72 & \tt B900 & \tt FDB8 & \tt FB34 & \tt FF38 & \tt 991E \\
\hline
\end{tabular}
\end{center}
\medskip

The basis satisfies $\beta_0 = 1$ and $\beta_i^2 + \beta_i = \beta_{i-1}$ for
$1 \le i \le 15$, which is a useful check for an implementation.
A stored 16-bit symbol value $x = \sum_b x_b 2^b$ denotes the field element

\[ \langle x \rangle = \sum_{b=0}^{15} x_b\,\beta_b . \]

Addition of symbols is therefore a bitwise XOR of their values, while
multiplication has to be carried out on $\langle\cdot\rangle$. Leopard-RS
does this with logarithm tables taken with respect to the generator $z$
({\tt 0x0002} in polynomial representation), indexed directly by the
stored symbol value. Examples (right hand sides in polynomial
representation): $\langle 1 \rangle = {\tt 0001}$,
$\langle 2 \rangle = {\tt ACCA}$, $\langle 3 \rangle = {\tt ACCB}$,
$\langle 4 \rangle = {\tt 3C0E}$, and
$\langle 2 \rangle \cdot \langle 3 \rangle = \langle 1 \rangle$.

\subsubsection{Codewords}
\label{rs04-codeword}

Let $k \ge 1$, $1 \le m \le k$, $M = {\rm NextPow2}(m)$ and $M + k \le 65536$.
The evaluation points are $\omega_j = \langle j \rangle$ for
$0 \le j < M + k$; they are pairwise different. An RS04 codeword with the
original symbols $d_0, \ldots, d_{k-1}$ is the unique vector
$(c_0, \ldots, c_{M+k-1})$ of symbols with

\[ c_{M+i} = d_i \quad (0 \le i < k) \qquad \mbox{and} \qquad
   \sum_{j=0}^{M+k-1} \langle c_j \rangle \, \omega_j^{\,t} = 0
   \quad (0 \le t < M). \qquad (1)\]

For given $d_i$, (1) is a linear system for $c_0, \ldots, c_{M-1}$ whose
matrix $(\omega_j^{\,t})$ is a Vandermonde matrix of distinct points, so
the solution exists and is unique. The {\em parity symbols} are
$c_0, \ldots, c_{m-1}$; $c_j$ is stored in ecc layer $j$. The symbols
$c_m, \ldots, c_{M-1}$ are never stored.

\medskip

The equations (1) define a generalized Reed-Solomon code of length $M+k$ and
dimension $k$ (the dual of the code of polynomials of degree $< M$ evaluated
at $\omega_j$). It is MDS, and so is the stored code of length $n = k + m$
obtained by omitting $c_m, \ldots, c_{M-1}$: {\bf any $k$ of the $n$ stored
symbols determine the codeword}, i.e. up to $m$ erasures per codeword can
be corrected.

\paragraph{Special cases.} For $m = 1$ ($M = 1$), (1) reduces to
$c_0 = d_0 \oplus d_1 \oplus \cdots \oplus d_{k-1}$ (XOR parity). For
$k = 1$, $m$ must be 1 and the parity is a copy of the original. Both
cases need no special treatment; they follow from (1).

\paragraph{Limits.} The limit is $M + k \le 65536$ and {\em not}
$n \le 65536$: $k = 49152, m = 16384$ is valid, but
$k = 49153, m = 16383$ is not. $k \le 65535$ follows.

\paragraph{Shards.} The code is applied to shards symbol by symbol: for
every symbol index $s$ ($0 \le s < B_s/2$), symbol $s$ of the $k$ original
shards and of the $m$ parity shards of a codeword satisfy (1).

\subsubsection{Computing the code (informative)}

The code (1) is exactly the GF($2^{16}$) code of Leopard-RS by Christopher
A.\ Taylor (\url{https://github.com/catid/leopard}, file
{\tt LeopardFF16.cpp} of the commit of 2024-03-23

\centerline{\tt 6e5725ebdf9da4370b0bcc4f70fa8eb66f4e6198}

which is the reference for bit-exactness). Leopard computes it with the
additive FFT over the novel polynomial basis of Lin, Chung and Han:

\begin{itemize}
\item S.-J. Lin, W.-H. Chung, Y. S. Han: {\em Novel polynomial basis and its
  application to Reed-Solomon erasure codes}, FOCS 2014.
\item S.-J. Lin, T. Y. Al-Naffouri, Y. S. Han, W.-H. Chung: {\em Novel
  polynomial basis with fast Fourier transform and its application to
  Reed-Solomon erasure codes}, IEEE Trans. Inf. Theory 62(11), 2016.
\end{itemize}

The encoder takes the originals in chunks of $M$, computes the inverse FFT
of each chunk with skew offset $M(i+1)$ for chunk $i$ (zero-padding the last
chunk), XORs the results and applies a forward FFT of size $M$; output $j$
is $c_j$. The decoder evaluates an error locator with a Walsh-Hadamard
transform modulo 65535 and needs work space for
$N = {\rm NextPow2}(M + k)$ shards (the {\em FFT size}). A skew factor
logarithm of 65535 means ``multiply by zero'' (skip). The log/exp tables,
skew factors and Walsh-Hadamard table are built exactly as in
{\tt LeopardFF16.cpp}.

\medskip

Usable implementations are {\tt src/rs04-ff16.c} (a C port of the FF16 code
path), Leopard-RS built with GF($2^{16}$) only, and the high-rate engine of
the Rust crate {\tt reed-solomon-simd}. {\bf Caution:} the stock
{\tt leo\_encode()} switches to GF($2^8$) when
${\rm NextPow2}(M+k) \le 256$ and then produces different parity (except for
$m = 1$ and $k = 1$); RS04 always uses GF($2^{16}$). Any other method, e.g.
solving (1) directly, gives the same result. {\tt contrib/rs04-reference.py}
is a slow reference which does exactly that.

\paragraph{Example.} For a single symbol column with $k = 3$, $m = 2$ and
originals {\tt 0001 0002 0003}, the parity symbols are {\tt 000C 000C}.
With $k = 5$, $m = 3$ ($M = 4$) and originals {\tt 0001}\ldots{\tt 0005}
they are {\tt 000E 0001 000C}. With $m = 1$ the parity of
{\tt 0001 0002 0003} is {\tt 0000}.

\subsection{Layout of an RS04 set}
\label{rs04-layout}

\subsubsection{Codewords}

The {\em original sectors} $X_0, \ldots, X_{kL-1}$ are the data sectors,
followed by the CRC sectors, followed by zero sectors:

\[ X_t = \left\{ \begin{array}{ll}
  D_t & 0 \le t < S \\
  C_{t-S} & S \le t < D \\
  \mbox{all-zero sector} & D \le t < kL
  \end{array} \right. \]

The zero sectors are virtual: they are never stored and never read, and
decoders treat them as known (not as erasures). $kL - D$ may exceed $L-1$,
i.e. complete {\em padding layers} are allowed; they make redundancies above
100\% possible in spite of $m \le k$.

\medskip

{\bf Codeword $i$} ($0 \le i < L$) has the originals
$d_j = X_{jL+i}$ ($0 \le j < k$) and the parity sectors
$\Pi_{jL+i} = c_j$ ($0 \le j < m$), computed from the shards of these
sectors as described in sections \ref{rs04-shard} and \ref{rs04-code}.
Thus data layer $j$ consists of the sectors $jL \ldots jL+L-1$ of the data,
and ecc layer $j$ of the parity sectors $\Pi_{jL} \ldots \Pi_{jL+L-1}$.
This is the RS03 rule with a larger number of layers; there is no
permutation inside layers.

\subsubsection{The added zone}
\label{rs04-zone}

The added zone stores the $Q = c + mL$ {\em payload sectors}
$C_0, \ldots, C_{c-1}, \Pi_0, \ldots, \Pi_{mL-1}$ (in this order, payload
index $q = 0 \ldots Q-1$) and the tag sectors. Sector $a$ of the added zone
($0 \le a < A$) is

\begin{itemize}
\item tag sector number $g = a / T$ if $a \bmod T = 0$;
\item otherwise payload sector $q = a - \lfloor a/T \rfloor - 1$.
\end{itemize}

Conversely, payload sector $q$ is stored at
$a = q + \lfloor q/(T-1) \rfloor + 1$. There are $\lceil Q/(T-1) \rceil$
tags and $A = Q + \lceil Q/(T-1) \rceil$; the last group may be shorter
than $T$ sectors, but every tag is followed by at least one payload sector.
Sector $a$ of the added zone has the {\em relative position} $S + a$.

\medskip

Consequences: the first sector of the added zone, directly behind the data,
is a tag; any $T$ consecutive sectors of the added zone contain a tag; the
last $T$ sectors of the added zone contain a tag.

\subsubsection{Augmented images and ecc files}

An {\bf augmented image} consists of the $S$ data sectors followed by the
$A$ sectors of the added zone; it is $S + A$ sectors long. The data sectors
are unchanged, so the image remains usable (e.g. mountable).

\medskip

An {\bf ecc file} is the added zone alone: $A \cdot B$ bytes, sector $a$ at
byte offset $a \cdot B$. Its contents are identical to the added zone of the
augmented image, so concatenating the data and the ecc file yields the
augmented image and vice versa. There is no separate ecc header; the first
sector of the file is a tag. A tag begins with the same 16 bytes
{\tt *dvdisaster*RS04} that an {\tt EccHeader} with method ``RS04'' would
have, so programs which look for the ecc header cookie at the start of a
file recognise the method. All positions stored in tags are relative to the
set start and therefore the same in both cases.

\subsubsection{Example}

The format test vector F1 (section \ref{rs04-vectors}) has $B = 2048$,
$dataBytes = 81820$ (so $S = 40$, the last sector holds 1948 bytes), $c = 1$,
$D = 41$, $L = 7$, $k = 6$, $m = 3$ and $T = 8$. The originals are
$X_0 \ldots X_{39} = D_0 \ldots D_{39}$, $X_{40} = C_0$ and the zero sector
$X_{41}$. $Q = 22$, there are 4 tags and $A = 26$:

\begin{center}
\begin{tabular}{|l|l|}
\hline
added zone sector $a$ (relative position $40+a$) & contents \\
\hline
0 & tag 0 (covers $C_0, \Pi_0 \ldots \Pi_5$) \\
1 & $C_0$ \\
2 \ldots 7 & $\Pi_0 \ldots \Pi_5$ \\
8 & tag 1 (covers $\Pi_6 \ldots \Pi_{12}$) \\
9 \ldots 15 & $\Pi_6 \ldots \Pi_{12}$ \\
16 & tag 2 (covers $\Pi_{13} \ldots \Pi_{19}$) \\
17 \ldots 23 & $\Pi_{13} \ldots \Pi_{19}$ \\
24 & tag 3 (covers $\Pi_{20}$) \\
25 & $\Pi_{20}$ \\
\hline
\end{tabular}
\end{center}

Codeword 3 consists of the data sectors 3, 10, 17, 24, 31, 38 and the parity
sectors $\Pi_3, \Pi_{10}, \Pi_{17}$ at the relative positions 45, 53 and 61.
Codeword 5 contains $X_{40} = C_0$, and codeword 6 contains the zero sector
$X_{41}$ as its last original.

\subsection{CRC sectors}
\label{rs04-crc}

CRC sector $C_j$ ($0 \le j < c$) holds 512 little-endian 32-bit values;
value $u$ (bytes $4u \ldots 4u+3$) is the Crc32 of data sector
$D_{512j+u}$ if $512j + u < S$ and zero otherwise. The Crc32 is taken over
the $B$ bytes of the sector as it is encoded, i.e.\ including the zero
padding of the last data sector. Bytes $2048 \ldots B-1$ of a CRC sector are
zero.

\medskip

CRC sectors are originals of the code, so lost CRC sectors are recovered by
the error correction like data sectors (as with the RS03 CRC layer). They
have no header because every tag describes the layout, and their integrity
is protected by the Crc32 values in the tags.

\subsection{Tag sectors}
\label{rs04-tags}

The byte layout of a tag is given in appendix \ref{rs04-tag}. A tag holds
the complete layout of the set and the Crc32 of the payload sectors that
follow it: tag $g$ covers the payload sectors
$g(T-1) \le q < \min(Q, (g+1)(T-1))$ and stores the Crc32 of payload sector
$q$ in $crc[q - g(T-1)]$; unused entries are zero. Tags are not part of
the code. They are protected by replication instead: all tags of a set are
identical except for $ownPosition$, $selfCRC$ and $crc[\,]$.

\medskip

A sector is a {\bf valid tag} if all of the following holds; otherwise it
is not a tag (it may be a data or parity sector which looks similar by
chance):

\begin{enumerate}
\item bytes 0--15 are {\tt *dvdisaster*RS04};
\item $selfCRC$ equals the Crc32 of bytes 0--2047 of the sector with bytes
  28--31 set to zero;
\item $sectorBytes \in \{2048, 2324, 2352\}$ and $tagInterval$ is a power
  of two with $2 \le T \le 256$;
\item $dataBytes \ge 1$, $L \ge 1$, $1 \le m \le k$,
  ${\rm NextPow2}(m) + k \le 65536$ and $kL \ge D$;
\item $ownPosition \ge S$, $(ownPosition - S) \bmod T = 0$ and
  $ownPosition - S < A$.
\end{enumerate}

If 1 and 2 hold but $version$ or $layoutVariant$ is not supported, the tag
belongs to a set of a newer format; see section \ref{rs04-versions}.
Readers should use the values held by the majority of the valid tags of a
set if tags disagree.

\subsection{Checksums}
\label{rs04-checksums}

\begin{tabular}{|p{2.8cm}|p{2.2cm}|p{4.5cm}|p{4.5cm}|}
\hline
Checksum & Stored in & Covers & Purpose \\
\hline
Crc32 per data sector & CRC sectors & the $B$ bytes of $D_t$ (last sector
zero-padded) & detect damaged data sectors (erasures); audio read offset
detection \\
\hline
Crc32 per payload sector & tags & the $B$ bytes of each CRC and parity
sector & detect damaged CRC and parity sectors (erasures) \\
\hline
$selfCRC$ & each tag & bytes 0--2047 of the tag with $selfCRC = 0$ &
recognise valid tags \\
\hline
$dataMD5$ & each tag & exactly the $dataBytes$ bytes of data, without
padding & end-to-end check after repair; matching image and ecc file \\
\hline
\end{tabular}
\medskip

Zero sectors and the parity symbols $c_m \ldots c_{M-1}$ are not stored and
have no checksum. Unlike RS01/RS02 there is no MD5 sum over the
error correction data: the per-sector Crc32 values in the tags check every
CRC and parity sector individually and also tell {\em which} sector is
damaged, which a single MD5 sum cannot.

\subsection{Choosing the encoding parameters (recommended)}
\label{rs04-params}

Any parameters which satisfy the validity rules of section \ref{rs04-tags}
can be decoded. dvdisaster chooses them as follows.

\paragraph{1. Parity budget $E$.}
\begin{itemize}
\item {\em Ecc files} and explicitly requested redundancy:
  $E = \lfloor \rho D + 1/2 \rfloor$ (rounded half up) for the requested
  redundancy $\rho$ (parity sectors per protected sector). The size of the
  set is $C_{tot} = D + E$. dvdisaster takes $\rho$ from the option
  {\tt -n}: $p${\tt \%} gives $\rho = p/100$, $r${\tt r} (``roots'' as for
  RS03; also without the suffix) $\rho = r/(255-r)$, {\tt normal} (the
  default) $32/223$ and {\tt high} $64/191$. $x${\tt m} chooses the
  largest $\rho$ for which the ecc file has at most $x$ MiB.
\item {\em Augmented images (default): fill the medium.} Let $C$ be the
  number of sectors available for the whole set: the capacity of the target
  medium (as reported by the drive, from the medium size table also used by
  RS02/RS03, or given by the user), or for a further session of a
  multi-session disc the free space of that session. With $Z = C - S$ and
  $T_0 = {\rm TagInterval}(Z)$ the largest number of payload sectors that
  fits is $Q_{max} = Z - \lceil Z / T_0 \rceil$, and
  $E = Q_{max} - c$. $C_{tot} = C$.
\item In both cases $E$ is capped at $2D$ (200\% redundancy, as RS03).
  Encoding is refused if $E < 1$ or $E/(D+E) < 8/255$ (less than the 8
  roots RS03 requires as a minimum); a warning is given if $E < D/5$
  (less than 20\% redundancy).
\end{itemize}

\paragraph{2. FFT size limit $N$.}
\[ N = \min\left(65536,\ {\rm Pow2Floor}\left(\left\lfloor
    \frac{C_{tot}}{32} \right\rfloor\right),\ {\rm Pow2Floor}\left(
    \left\lfloor\frac{R}{128 \cdot B_s}\right\rfloor\right)\right) \]
where $R$ is the memory budget of the repair (default 4 GiB, dvdisaster
option {\tt --rs04-memory}). A repair pass
holds 128 consecutive codewords of up to $N$ shards each in memory (with
fewer codewords per pass a hard disk seeks too much); with 4 GiB this gives
$N = 16384$ for DVD and BD. The $C_{tot}/32$ term keeps $L \ge 32$ on small media
($N = 8192$ for CD). The user may override $N$. If no geometry satisfies
step 3 with this $N$ (only for very small sets), $N = 65536$ is used.

\paragraph{3. Geometry.} $L$ is the smallest integer
$L \ge \min(32, D, E)$ for which
\[ m = \lfloor E / L \rfloor \ge 1, \quad
   k = \max(\lceil D/L \rceil, m) \quad \mbox{and} \quad
   {\rm NextPow2}(m) + k \le N . \]
The left side of the last condition does not increase with $L$, so $L$ can
be found by bisection. If no $L \le E$ qualifies, encoding is refused.
$k > \lceil D/L \rceil$ (padding layers) only happens for redundancies above
100\%. The smallest $L$ gives the longest codewords for the memory budget,
and at most $L-1$ sectors of the budget remain unused.

\paragraph{4. Tag interval.} $T = {\rm TagInterval}(Q)$ with $Q = c + mL$,
where
\[ {\rm TagInterval}(x) = \left\{ \begin{array}{ll} 256 & x \ge 4080 \\
   \max(2, {\rm Pow2Floor}(\lfloor x/16 \rfloor)) & \mbox{otherwise}
   \end{array}\right. \]
This gives every set at least 16 tags (or $Q$ tags if $Q < 32$); large sets
spend $1/256$ of the added zone on tags. When filling a medium, $T = T_0$ is
used instead if the set would not fit with $T$.

\paragraph{5. Set ID.} 8 random bytes, newly chosen for every encoding run.

\medskip

{\bf Rationale for $L \ge 32$:} DVD drives lose data in units of 16 sectors
(one ECC block) and BD drives in units of 32 sectors (one cluster). With
$L \ge 32$ such a unit costs each codeword at most one symbol.

\medskip

{\bf Examples} (produced by {\tt contrib/rs04-reference.py}; ``fill'' uses
the medium sizes of table \ref{layout-size-table}):

\medskip
{\footnotesize
\begin{tabular}{|l|r|r|r|r|r|r|r|r|}
\hline
Case & $S$ & $L$ & $k$ & $m$ & $T$ & FFT size & set size & parity/$D$ \\
\hline
CD, fill & 300000 & 49 & 6135 & 1196 & 256 & 8192 & 359423 & 19.5\% \\
DVD, fill & 2000000 & 142 & 14113 & 2042 & 256 & 16384 & 2295024 & 14.5\% \\
DVD-9, fill & 3500000 & 286 & 12262 & 2315 & 256 & 16384 & 4171550 & 18.9\% \\
BD-25, fill & 10000000 & 816 & 12279 & 2205 & 256 & 16384 & 11825945 & 18.0\% \\
CD, fill, tiny image & 100 & 32 & 6 & 6 & 8 & 16 & 321 & 190\% \\
ecc file, $\rho = 32/223$ & 2000000 & 141 & 14213 & 2039 & 256 & 16384 & 2292549 & 14.3\% \\
CD-DA ecc file, 20\% & 330000 & 54 & 6124 & 1224 & 256 & 8192 & 397003 & 20.0\% \\
ecc file, 50\%, 10 sectors & 10 & 6 & 2 & 1 & 2 & 4 & 24 & 54.5\% \\
\hline
\end{tabular}
}

\subsection{Multi-session discs}
\label{rs04-sessions}

Each session of a multi-session disc may carry its own RS04 set, which
protects the sectors of that session only. The set is created from the
session image (as written by the burning software, e.g. the image created
for {\tt growisofs -M}) without knowing its absolute position on the disc:
all positions in the set are relative to its first data sector. When the
whole disc is read, the absolute start of a set is

\[ start = (\mbox{absolute LBA of a tag}) - ownPosition , \]

and the set occupies the absolute sectors $start \ldots start + S + A - 1$.
Gaps between sessions (lead-out, lead-in) belong to no set. An ISO 9660
file system in a later session may refer to files of earlier sessions;
those sectors are protected by the set of the session they are stored in.

\medskip

{\em CD Extra} (an audio session followed by a data session) is covered by
the same rules: the audio session is protected by an ecc file with
$B = 2352$ (section \ref{rs04-audio}), which may be stored as a file in the
data session, and the data session may carry its own augmented set.

\subsection{Recognising RS04 data}
\label{rs04-recognize}

\paragraph{Augmented images and media.} A reader looks for valid tags in
the following order and stops as soon as it has found what it needs:

\begin{enumerate}
\item If the image (or, on a multi-session medium, a session) starts with an
  ISO 9660 file system of $V$ sectors, the sectors $V$ and $V+150$ relative
  to its start (behind the file system, and behind the 150 sectors of
  padding added by some mastering programs). For an image which was not
  padded after the file system, sector $V$ is tag 0.
\item The last 256 sectors of the image or session. They contain a tag if
  the image ends with the complete added zone (since $T \le 256$).
\item All sectors, from the end towards the start (the added zone is at the
  end of a set).
\end{enumerate}

Every valid tag yields a {\em candidate set} $(setID, start)$ with
$start = position - ownPosition$; candidates with $start < 0$ are
discarded. Tags with equal $setID$ and $start$ belong to the same set. A
set may extend beyond the end of the image (a truncated image is repaired by
recovering the missing sectors). The size of the complete image is
$start + S + A$; dvdisaster uses it as the expected image size when reading
a medium.

\medskip

To enumerate all sets of a multi-session disc, the reader applies steps 1
and 2 to every session in the table of contents, or step 3 to a disc image
without table of contents, and reports every candidate set.

\medskip

\paragraph{Stale and embedded sets.} Candidates may also come from tags of
an earlier recording on rewritable media, or from an ecc file or augmented
image which is stored as a {\em file} inside the protected data. A reader
confirms a candidate by comparing up to 16 readable data sectors with their
Crc32 values from readable CRC sectors. A candidate whose compared sectors
all mismatch is rejected. If confirmed candidates overlap, the one with more
matching sectors wins, then the one with more valid tags. A candidate
without readable data or CRC sectors is kept but reported as unconfirmed.

{\em Informative:} dvdisaster applies steps 1 and 2 to every image it
opens. Step 3 is applied to the parts of an image file which are not covered
by the sets found so far (gaps between and before sessions), and to a whole
image without any set found only on request (option {\tt -a RS04}), since
reading every sector of every opened image would be too slow. A set whose
tag 0 is lost and whose image has more than 256 sectors appended is
therefore only found with {\tt -a RS04}. On a medium, steps 1 and 2 are used.

\paragraph{Ecc files.} Sector 0 of an ecc file is tag 0. If it is
unreadable or invalid, the reader searches the file for the 16-byte
signature at byte offsets which are multiples of 4 (all three sector sizes
are multiples of 4) and accepts a valid tag at offset $o$ only if
$o = (ownPosition - S) \cdot sectorBytes$. The ecc file must be
$A \cdot B$ bytes long; shorter files are treated as truncated (the missing
sectors are erasures).

\medskip

Recognition never depends on the medium size, the medium type or a table of
known sizes.

\subsection{Audio CDs and read offsets}
\label{rs04-audio}

Audio CDs are protected by ecc files made from a rip which covers a fixed
range of the disc (normally all tracks as one {\tt .bin} file, $B = 2352$).
CD drives differ in their {\em read offset}: two rips of the same disc made
with different drives (or different offset corrections) are shifted by a
whole number of samples (4 bytes) against each other. RS04 finds the shift
with the data sector Crc32 values.

\medskip

Let $R$ be the new rip (a byte string) and $crc_t$ the stored Crc32 of data
sector $t$. For an offset $\delta$ (in samples, may be negative) the
{\em aligned sector} $t$ is the 2352 bytes of $R$ starting at byte
$2352\,t + 4\delta$. The recommended procedure is:

\begin{enumerate}
\item {\bf Probe sectors.} Take every data sector $t$ with
  $10 \le t < S - 10$ whose $crc_t$ is known (from a readable or recovered
  CRC sector) and differs from {\tt 0x416831C0}, the Crc32 of digital
  silence (silence matches at every offset). Choose up to 64 of them, evenly
  spread over this range.
\item {\bf Fast path.} If all probe sectors match for $\delta = 0$, the rip
  is aligned.
\item {\bf Search.} For every $\delta$ with $|\delta| \le 5880$ samples
  (10 sectors; drive read offsets are usually below 2000 samples in
  magnitude, and two rips differ by the difference of their offsets), the score is the
  number of probe sectors whose aligned sector has the Crc32 $crc_t$. Probe
  sectors whose aligned sector is not completely inside $R$ or touches a
  sector which was unreadable in the new rip do not count.
\item {\bf Decision.} The offset with the highest score is accepted if the
  score is at least 2 and at least twice the second highest score. Otherwise
  (e.g. too much damage, or the audio repeats exactly with a period of whole
  samples) the user has to specify the offset (dvdisaster option
  {\tt --audio-offset}); manual offsets are not limited to the search window.
\end{enumerate}

With the accepted $\delta$, aligned sectors which are not completely inside
$R$ or overlap unreadable parts of the rip are erasures. With $\delta \ne 0$
this concerns the first or the last $\lceil 4|\delta| / 2352 \rceil$ sectors,
because drives normally cannot read into the lead-in or lead-out; they are
recovered like any other erasure, and the repaired rip is identical to the
original one ($dataMD5$ matches). dvdisaster repairs such a rip by first
rewriting the file shifted by $4\delta$ bytes (missing bytes at the edges
become zero) and then decoding it with $\delta = 0$. Rips which differ in more than a constant
shift (other track splitting, pregaps or hidden track audio removed) are out
of scope; the user has to provide a rip of the same range. Reading audio
sectors from a drive (READ CD with C2 error pointers) does not affect the
format.

\subsection{Decoding (informative)}
\label{rs04-decode}

\begin{enumerate}
\item Recognise the set and read its tags (ignoring tags with another
  $setID$), CRC sectors and parity sectors.
  An added zone sector is an erasure if it is unreadable or its Crc32 does not
  match the value in its tag. If its tag is lost, a readable sector is
  assumed to be correct.
\item If CRC sectors are missing, decode the codewords which contain them
  first ($C_j$ is in codeword $(S+j) \bmod L$). Data sectors whose Crc32 is
  not known yet are assumed to be correct if they are readable.
\item For every codeword: a data sector is an erasure if it is unreadable
  (in image files: contains the dead sector marker of dvdisaster, section
  \ref{rs04-missing}) or its
  Crc32 does not match. Zero sectors are known. With at most $m$
  erasures the codeword is decoded, otherwise its erased sectors stay
  unrecoverable. The data sector Crc32 values are checked after decoding.
\item After a complete repair, the MD5 sum of the data must equal
  $dataMD5$.
\end{enumerate}

The decoder only needs one $B_s$-byte shard per layer and codeword at a
time; an implementation processes $W$ consecutive codewords per pass to read
the image in larger pieces.

\subsection{Versioning and compatibility}
\label{rs04-versions}

\begin{itemize}
\item $version$ identifies the format of tags, CRC sectors and the code.
  This specification defines version 1. A reader must refuse to process a
  set whose tags carry a different version and should tell the user that a
  newer dvdisaster is required. Incompatible changes increase the version.
\item $layoutVariant$ identifies the mapping of codewords to sectors.
  0 is the layout of this specification (contiguous layers, no
  permutation). Readers must refuse sets with an unknown layout variant.
\item $sectorBytes$ values other than 2048, 2324 and 2352 must be refused.
\item Unknown bits in $flags$ are ignored; flags never change how data
  is decoded.
\item Reserved bytes (the unused tail of the tag and of CRC sectors) are
  written as zero and ignored when read.
\item $creatorVersion$ is informational (diagnostics and workarounds for
  bugs of a particular dvdisaster version).
\end{itemize}

\subsection{Limits and edge cases}
\label{rs04-limits}

\begin{tabular}{|p{4.2cm}|p{10.3cm}|}
\hline
Data size & $1 \le dataBytes < 2^{64}$. An empty image cannot be protected.
Augmented images need $dataBytes = S \cdot B$. \\
\hline
Codeword & $1 \le m \le k$ and ${\rm NextPow2}(m) + k \le 65536$, hence
$k \le 65535$ and $n \le 65536$; the FFT size is at most 65536. \\
\hline
Layer size & $L \ge 1$ and $kL \ge D$. There is no upper limit; with
$k = 1$ a layer holds all protected sectors. \\
\hline
Small sets & No special cases. $L < 32$ is allowed (dvdisaster only uses it
when $D < 32$ or $E < 32$); then one 32-sector unit may hit a codeword
twice. A minimal set: $S = 1$, $L = 1$, $k = 2$, $m = 1$, $T = 2$,
$Q = 2$, $A = 4$. \\
\hline
$k = 1$ & $m = 1$; each parity sector is a copy of the original sector. \\
\hline
$m = 1$ & XOR parity over the $k$ originals. \\
\hline
Zero sectors & Up to $kL - D$ sectors (possibly whole layers) are virtual
zeros; they are never stored. \\
\hline
Tags & At least one ($Q \ge 2$). The last group of the added zone may
have fewer than $T-1$ payload sectors. \\
\hline
Implementation & {\em Informative:} dvdisaster handles sets with
$S$, $kL$ and $A$ below $2^{31}$ sectors (4 TiB of 2048-byte sectors) and
treats tags describing larger sets as invalid. \\
\hline
\end{tabular}

\subsection{Design rationale}
\label{rs04-rationale}

\paragraph{Damage simulation.} The layout was chosen with the damage
simulator {\tt contrib/rs04-layout-sim.py} (Python 3 and numpy):

{\small
\begin{verbatim}
python3 contrib/rs04-layout-sim.py --trials 200 > results.md
\end{verbatim}
}

({\tt --quick} gives a short smoke run). It maps damage models onto sectors using the
CLV geometry of CD-R, DVD and BD (outer edge bands on single and dual layer
discs, scattered loss of 16/32-sector ECC units, band with decaying fringe,
radial scratches, contiguous runs of 4096 sectors and 1~cm$^2$ spots) and
counts the erasures per codeword for RS03 and for RS04 layouts with
$n \approx 255 \ldots 65535$, with and without permutations inside layers.
The following table shows the loss (in percent of all sectors) at which the
probability of a complete recovery drops below 99\%, for a parity share of
$m/n = 32/255$ (MDS limit 12.5\%):

\medskip
{\small
\begin{tabular}{|l|r|r|r|r|r|}
\hline
Damage model, medium & RS03 & $n=1023$ & $n=4095$ & $n=16383$ & $n=65535$ \\
\hline
outer band, BD-25 & 12.55 & 12.51 & 12.55 & 12.55 & 12.59 \\
band, DVD-9 (opposite track path) & 12.55 & 12.51 & 12.56 & 12.57 & 12.62 \\
scattered units, CD-R 80 min & 5.46 & 8.50 & 10.69 & -- & -- \\
scattered units, DVD-5 & 5.27 & 8.61 & 10.59 & 11.78 & -- \\
scattered units, BD-25 & 4.72 & 8.48 & 10.49 & 11.59 & 12.14 \\
band with fringe, BD-25 & 8.16 & 10.29 & 11.42 & 12.10 & 12.40 \\
radial scratches, BD-25 & 5.76 & 8.80 & 10.85 & 11.62 & 12.25 \\
runs of 4096 sectors, BD-25 & 5.87 & 9.82 & 11.84 & 12.34 & 12.56 \\
spots, BD-25 & 6.11 & 9.21 & 10.96 & 11.77 & 12.26 \\
\hline
\end{tabular}
}
\medskip

Findings: (1) long codewords more than double the tolerable scattered loss
and stay at the limit for edge damage; (2) permuting 32-sector blocks inside
the layers is within the noise for $n \ge 4095$, worse for small $n$ and
for contiguous runs, and costs about $C/32$ seeks per repair, so RS04 uses
no permutation and only reserves the $layoutVariant$ field; (3) repair time
on a hard disk grows with $n$ for a given memory budget (BD-25 with 4 GiB:
about 4 minutes for $n \approx 1023$, 19 minutes for $n \approx 16323$ and
1.1 hours for $n \approx 65283$; about 1--2 minutes on an SSD), therefore
$n$ is chosen per image and stored in the tags; (4) dvdisaster issue \#43
is not explained by aliasing between the RS03 stride and the disc
revolution.

\paragraph{Other decisions.}
\begin{itemize}
\item {\em Tags in the added zone at a fixed interval} instead of a header
  at one place: any surviving part of the added zone describes the layout,
  also on dual layer discs whose outer edge is in the middle of the LBA
  range. Tag 0 is where RS03 has its ecc header, directly behind the data.
\item {\em Payload Crc32 values in the tags}: they cannot be in the CRC
  sectors, which are protected by the parity they would describe. The tags
  have room for them anyway.
\item {\em CRC sectors as originals} (like the RS03 CRC layer): lost
  checksums are recovered by the code. They have no header since the tags
  carry the layout.
\item {\em Virtual zero sectors and shard padding}: nothing is stored that
  can be computed, which makes the ecc file identical to the added zone.
\item {\em Relative positions and a random set ID}: sessions can be
  augmented without knowing where they will be written, and tags of other
  or older sets are never mixed in.
\item {\em One field for the data length} ($dataBytes$) instead of a sector
  count and the bytes in the last sector: both follow from it.
\item {\em No medium fingerprint}: the Crc32 values of the data sectors in
  the CRC sectors identify the data more reliably than the MD5 sum of one
  sector.
\end{itemize}

\subsection{Test vectors}
\label{rs04-vectors}

\paragraph{Code.} The originals of each vector are filled with the output
of the generator splitmix64 started with $state = seed$ (each step:
$state \leftarrow state + {\tt 0x9E3779B97F4A7C15}$, $z \leftarrow state$,
$z \leftarrow (z \oplus (z \gg 30)) \cdot {\tt 0xBF58476D1CE4E5B9}$,
$z \leftarrow (z \oplus (z \gg 27)) \cdot {\tt 0x94D049BB133111EB}$,
output $z \oplus (z \gg 31)$, all modulo $2^{64}$), each 64-bit output
stored little-endian, filling original shards $0 \ldots k-1$ consecutively.
The table gives the MD5 sum of the parity shards $0 \ldots m-1$
concatenated (shards in Leopard buffer format, section \ref{rs04-shard}).

\medskip
{\footnotesize
\begin{tabular}{|r|r|r|r|l|l|}
\hline
seed & $k$ & $m$ & bytes & MD5 of parity & case \\
\hline
1 & 2 & 1 & 64 & {\tt 9b1fe76ef28bf3823a5e3cf97adf0901} & XOR \\
2 & 5 & 3 & 64 & {\tt 74130cbd984417137639d2a582b9b4a1} & small$^*$ \\
3 & 100 & 31 & 128 & {\tt 8c5370f6adebf2d28b46125fdd201dc8} & no powers of 2 \\
4 & 1000 & 1000 & 64 & {\tt 82c42efe91c5110bb60f9235e7bd3991} & $m = k$ \\
5 & 3584 & 512 & 2048 & {\tt 047ba4c7183fabee56efa1053807edae} & typical \\
6 & 49152 & 16384 & 64 & {\tt 055b961a368a2f12c3b5526636eeaf31} & $M + k = 65536$ \\
\hline
\end{tabular}
}
\medskip

$^*$ stock {\tt leo\_encode()} would use GF($2^8$) for this size.
``bytes'' is the shard size. These values were produced by {\tt src/rs04-ff16.c}, whose output was
compared with Leopard-RS (built with GF($2^{16}$) only) and the high-rate
encoder of {\tt reed-solomon-simd}. Vectors 1--3 are also reproduced by
{\tt contrib/rs04-reference.py} directly from equation (1).

\paragraph{Format.} The data of a format vector is the splitmix64 stream
started with $seed$, truncated to $dataBytes$ bytes. All tags use
$setID = {\tt 01\ 02\ 03\ 04\ 05\ 06\ 07\ 08}$, $creatorVersion = 0$ and
$flags = 0$. The table gives the MD5 sums of the data and of the complete
added zone (= the ecc file).

\medskip
{\footnotesize
\begin{tabular}{|l|r|r|r|r|r|r|r|r|r|r|r|}
\hline
 & seed & $B$ & $dataBytes$ & $S$ & $c$ & $L$ & $k$ & $m$ & $T$ & $Q$ & $A$ \\
\hline
F1 & 7 & 2048 & 81820 & 40 & 1 & 7 & 6 & 3 & 8 & 22 & 26 \\
F2 & 8 & 2352 & 23520 & 10 & 1 & 4 & 3 & 2 & 4 & 9 & 12 \\
F3 & 9 & 2324 & 11620 & 5 & 1 & 6 & 1 & 1 & 2 & 7 & 14 \\
\hline
\end{tabular}

\medskip
\begin{tabular}{|l|l|l|}
\hline
 & MD5 of the data & MD5 of the added zone \\
\hline
F1 & \tt ac8364a7c7b4fcac9126f9d120a2b243 & \tt ce7142e2d9767ac6642ea567df200f90 \\
F2 & \tt b5175b33aa5f95bec590eec94cb8392d & \tt 84d948d015532a395e1ccdea5e19b62f \\
F3 & \tt 225073ed5626b7225ec38cc288d1a6d2 & \tt 28a4b07832bda689aaa9b4e1278333ca \\
\hline
\end{tabular}
}
\medskip

F1 has a partial last data sector and a padding sector, F2 exercises the
sector to shard mapping, F3 the case $k = m = 1$ with $T = 2$. The first 96
bytes of tag 0 of F1 are:

{\footnotesize
\begin{verbatim}
0000: 2a 64 76 64 69 73 61 73 74 65 72 2a 52 53 30 34  *dvdisaster*RS04
0010: 01 00 00 00 00 08 08 00 00 00 00 00 ce bb 08 3b  ................
0020: 01 02 03 04 05 06 07 08 9c 3f 01 00 00 00 00 00  .........?......
0030: 07 00 00 00 00 00 00 00 06 00 00 00 03 00 00 00  ................
0040: 28 00 00 00 00 00 00 00 ac 83 64 a7 c7 b4 fc ac  (.........d.....
0050: 91 26 f9 d1 20 a2 b2 43 f1 75 b2 b2 c8 97 c7 92  .&.. ..C.u......
\end{verbatim}
}

\fussy
